\section{Covariance, forme normale et spectre gaussien}
\label{app:gaussianas-espectro}

La section positive d’action \(\hbar_a\) et la réalisation canonique
précédente déterminent la normalisation employée ici. On démontre la
correspondance entre covariance, spectre, pureté et entropie pour un
nombre fini de modes ; l’espace de Fock conserve toutes ses occupations.
La réduction symplectique et les formules gaussiennes sont des résultats
établis qui sont reproduits pour rendre autonome leur application au
transport de mémoire.

\subsection{Coordonnées canoniques et fonction caractéristique}
Dans \(L^2(\RR^d)\), soient \(\widehat q_j=\widehat Q_j/\sqrt{\hbar_a}\),
\(\widehat p_j=\widehat P_j/\sqrt{\hbar_a}\) et
\(\widehat z=(\widehat q_1,\widehat p_1,\ldots,\widehat q_d,\widehat p_d)\).
Sur Schwartz,

\[
 [\widehat z_j,\widehat z_k]=iJ_{jk}I,\qquad
 J=\bigoplus_{j=1}^d\begin{pmatrix}0&1\\-1&0\end{pmatrix}.
\]
Cette convention fixe le signe de la forme canonique. Les opérateurs de
Weyl \(\mathcal W(\xi)=\exp(i\xi^{\mathsf T}\widehat z)\) sont définis
par les fermetures auto-adjointes des combinaisons linéaires indiquées.
Pour un état centré \(\rho\) de moments d’ordre deux finis, la
covariance normalisée et sa fonction caractéristique sont
\[
 W_{jk}=\Tr\rho(\widehat z_j\widehat z_k+\widehat z_k\widehat z_j),
 \qquad \chi_\rho(\xi)=\Tr(\rho\mathcal W(\xi)),\qquad
 V=\frac{\hbar_a}{2}W.
\]
L’état est gaussien centré lorsque

\begin{equation}
 \chi_\rho(\xi)=\exp(-\xi^{\mathsf T}W\xi/4).
 \label{app:eq:caracteristica-gaussiana}
\end{equation}
L’espérance de \((c^{\mathsf T}\widehat z)^\dagger
(c^{\mathsf T}\widehat z)\) est \(c^\dagger(W+iJ)c/2\) ; ainsi,
\(W+iJ\geq0\). De plus, \(W>0\) : si \(v\) est réel et
\(v^{\mathsf T}Wv=0\), la positivité implique \((W+iJ)v=0\),
donc \(Jv=0\) et \(v=0\).


\begin{lemma}[Unicité de la fonction caractéristique
]
\label{app:lem:unicidad-weyl}
Deux opérateurs de classe trace ayant la même fonction caractéristique sont égaux.
\end{lemma}
\begin{proof}
En regroupant les positions et les moments sous la forme \(\xi=(u,v)\),

\[
 (\mathcal W(u,v)f)(x)=e^{iu\cdot(x+v/2)}f(x+v).
\]
Pour un opérateur de rang fini à vecteurs de Schwartz et de noyau \(K_A\),

\[
 \Tr(A\mathcal W(u,v))=\int K_A(x+v/2,x-v/2)e^{iu\cdot x}\,dx.
\]
Plancherel et le changement de variables de jacobien absolu un donnent
\[
 \|A\|_{\rm HS}^2=(2\pi)^{-d}\int_{\RR^{2d}}
 |\Tr(A\mathcal W(\xi))|^2\,d\xi.
\]
Les opérateurs précédents sont denses dans la classe trace. La convergence
dans cette norme implique la convergence en Hilbert--Schmidt et la convergence
uniforme des fonctions caractéristiques, car leur différence est
bornée par \(\|A-B\|_1\). L’égalité s’étend dans \(L^2\).
L’appliquer à la différence des opérateurs prouve l’assertion.
\end{proof}

\subsection{Réduction symplectique de la forme positive
}
\begin{theorem}[Forme normale de Williamson]
\label{app:thm:williamson}
Pour une matrice réelle symétrique \(W>0\) d’ordre \(2d\), il existe
\(S\in\operatorname{Sp}(2d,\RR)\) et \(\nu_j>0\) tels que
\[
 W=SDS^{\mathsf T},\qquad D=\bigoplus_j\nu_jI_2.
\]
Les valeurs propres de \(JW\) sont \(\pm i\nu_j\). La condition
\(W+iJ\geq0\) équivaut à \(\nu_j\geq1\) pour tout \(j\).

\end{theorem}
\begin{proof}
La matrice \(B=W^{1/2}JW^{1/2}\) est antisymétrique et inversible.
Si \(-B^2u=\nu^2u\) et \(\|u\|=1\), le vecteur
\(v=-Bu/\nu\) est unitaire et orthogonal à \(u\) ; il satisfait
\(Bu=-\nu v\), \(Bv=\nu u\). Le complément de cette paire est
invariant. L’itération produit une matrice orthogonale \(O\) telle que
\(O^{\mathsf T}BO=JD\). En définissant \(S=W^{1/2}OD^{-1/2}\),

\[
 S^{\mathsf T}JS=J,\qquad SDS^{\mathsf T}=W.
\]
La similitude \(W^{1/2}(JW)W^{-1/2}=B\) détermine les valeurs
\(\nu_j\), avec multiplicité. La congruence par \(S^{-1}\)
transforme \(W+iJ\) en \(D+iJ\) ; chaque bloc a pour valeurs propres
\(\nu_j+1\) et \(\nu_j-1\).
\end{proof}

\begin{lemma}[Implémentation unitaire dans un nombre fini de modes]
\label{app:lem:implementacion-simplectica}
Toute matrice symplectique \(S\) admet un unitaire \(U_S\) qui satisfait
\(U_S^\dagger\mathcal W(\xi)U_S=\mathcal W(S^{\mathsf T}\xi)\).
La conjugaison \(\rho\mapsto U_S\rho U_S^\dagger\) transporte
la covariance par \(W\mapsto SWS^{\mathsf T}\).
\end{lemma}
\begin{proof}
Dans la décomposition polaire \(S=OP\), l’identité
\(J^{-1}(S^{\mathsf T}S)J=(S^{\mathsf T}S)^{-1}\), appliquée par
calcul spectral à la racine positive, prouve \(PJP=J\).
Alors \(O\) est orthogonale symplectique. Les espaces propres de \(P\)
s’apparient au moyen de \(J\) avec les valeurs propres \(\lambda,\lambda^{-1}\) ;
celui de valeur propre un est \(J\)-invariant. Une matrice
orthogonale symplectique \(K\) réduit donc \(P\) à
\(\bigoplus_j\operatorname{diag}(e^{r_j},e^{-r_j})\).

La dilatation est mise en œuvre par

\[
 (\mathcal D_rf)(x)=e^{-\frac12\sum_jr_j}
 f(e^{-r_1}x_1,\ldots,e^{-r_d}x_d).
\]
Un changement de variables prouve son unitarité et la conjugaison de
\((\widehat q_j,\widehat p_j)\) en
\((e^{r_j}\widehat q_j,e^{-r_j}\widehat p_j)\).
Une transformation orthogonale symplectique commute avec \(J\) et
correspond à une matrice \(u\in U(d)\) sur
\(a_j=(\widehat q_j+i\widehat p_j)/\sqrt2\). Sa seconde
quantification \(\Gamma(u)=\bigoplus_{n\geq0}u^{\otimes_s n}\)
est unitaire dans \(\mathcal F_s(\CC^d)\), et l’identité sur les
produits symétriques finis donne la transformation des \(a_j\).

La base de Hermite identifie cet espace à \(L^2(\RR^d)\).
Pour justifier sa complétude, si \(f\) est orthogonal à tous les
polynômes multipliés par \(e^{-|x|^2/2}\), la fonction entière
\(F(z)=\int f(x)e^{-|x|^2/2}e^{z\cdot x}\,dx\) a toutes
ses dérivées nulles en zéro. Ainsi, \(F=0\) ; en restreignant à
\(z=i\xi\), l’unicité de Fourier donne \(f=0\). Les
normalisations et l’orthogonalité résultent de la création et de l’annihilation.
La composition des mises en œuvre précédentes produit \(U_S\).
L’égalité de Weyl implique
\(\chi_{U_S\rho U_S^\dagger}(\xi)=\chi_\rho(S^{\mathsf T}\xi)\),
qui démontre la loi de covariance.

\end{proof}

\subsection{Calcul du spectre de la densité isotrope}
Dans un mode, les états cohérents normalisés sont
\[
 |z\rangle=e^{-|z|^2/2}\sum_{n\geq0}\frac{z^n}{\sqrt{n!}}|n\rangle.
\]
Leur fonction de position,
\(\pi^{-1/4}\exp[-(x-q_0)^2/2+ip_0(x-q_0/2)]\), où
\(q_0=\sqrt2\Re z\), \(p_0=\sqrt2\Im z\), donne par intégration
\[
 \langle z|\mathcal W(u,v)|z\rangle
 =e^{-(u^2+v^2)/4}e^{i(uq_0+vp_0)}.
\]
Pour \(s>0\), l’intégrale
\[
 \rho^{(s)}=\frac1{\pi s}\int_{\CC} e^{-|z|^2/s}
 |z\rangle\langle z|\,d^2z
\]
converge en classe trace, est positive et a pour trace un. L’intégration
angulaire annule ses éléments hors de la diagonale d’occupation ; l’intégration
radiale donne
\[
 \langle n|\rho^{(s)}|n\rangle=\frac{s^n}{(1+s)^{n+1}},\qquad
 \chi_{\rho^{(s)}}(u,v)=e^{-(2s+1)(u^2+v^2)/4}.
\]
Avec \(\nu=2s+1\), l’unique état gaussien de covariance
\(\nu I_2\) est, par le lemme~\ref{app:lem:unicidad-weyl},
\begin{equation}
 \rho_\nu=(1-t)\sum_{n\geq0}t^n|n\rangle\langle n|,
 \qquad t=\frac{\nu-1}{\nu+1}.
 \label{app:eq:estado-geometrico}
\end{equation}
Le cas \(\nu=1\) est le projecteur sur le vide.

\begin{theorem}[Spectre, pureté et entropie
]
\label{app:thm:espectro-gaussiano}
Un état gaussien centré de covariance \(W\) et de valeurs
symplectiques \(\nu_1,\ldots,\nu_d\) est unitairement équivalent
à \(\bigotimes_j\rho_{\nu_j}\). Ses valeurs propres sont
\[
 \lambda_{\symbf{n}}=\prod_j(1-t_j)t_j^{n_j},\qquad
 t_j=\frac{\nu_j-1}{\nu_j+1},\quad \symbf{n}\in\mathbb N^d.
\]
Avec \(0\log0=0\),
\begin{align}
 \Tr\rho^2&=\prod_j\nu_j^{-1}=(\det W)^{-1/2},
 \label{app:eq:pureza-general}\\
 -\Tr(\rho\log\rho)&=\sum_j
 \left[\frac{\nu_j+1}{2}\log\frac{\nu_j+1}{2}
 -\frac{\nu_j-1}{2}\log\frac{\nu_j-1}{2}\right].
 \label{app:eq:entropia-general}
\end{align}
\end{theorem}
\begin{proof}
Williamson donne \(W=SDS^{\mathsf T}\). Le produit des densités
isotropes a pour fonction caractéristique \(e^{-\xi^{\mathsf T}D\xi/4}\).
L’implémentation unitaire et l’unicité de Weyl prouvent
l’équivalence et la formule spectrale. Dans un mode,
\[
 \sum_{n\geq0}(1-t)^2t^{2n}=\frac{1-t}{1+t}=\nu^{-1},
\]
et
\[
 -\sum_{n\geq0}(1-t)t^n\log((1-t)t^n)
 =-\log(1-t)-\frac{t}{1-t}\log t.
\]
La substitution \(t=(\nu-1)/(\nu+1)\) prouve les formules ; la
limite en \(t=0\) est nulle. Dans le produit fini, Tonelli justifie
la somme des termes d’entropie positifs ou nuls, et
\(\det W=\prod_j\nu_j^2\) donne l’expression déterminantale.

\end{proof}

Si l’état est la réduction d’un vecteur pur bipartite, ses
valeurs propres sont les carrés des coefficients de Schmidt.
La décomposition en valeurs singulières de l’opérateur de Hilbert--Schmidt défini
par ce vecteur démontre l’égalité des spectres non nuls des deux
réductions. L’entropie calculée est alors celle d’intrication.

\subsection{Préparation et transport des deux covariances}
L’application ultérieure utilise deux entrées gaussiennes pures identiques.
En général, pour \(M\in\operatorname{Sp}(2d,\RR)\), leur covariance est
\(V_0=(\hbar_a/2)MM^{\mathsf T}\), et l’entrée conjointe possède la
covariance \(V_0\oplus V_0\). Le transport physique s’obtient
en conjuguant \(U_H\) par \(M\oplus M\). Dans la coordinatisation équilibrée,
la covariance normalisée d’entrée est \(I_{4d}\) ; avec
\[
 T=\frac{8I+H}{9},\qquad D=\frac{\sqrt8(H-I)}9,
\]
la première réduction a
\[
 W=TT^{\mathsf T}+DD^{\mathsf T}=\frac{8I+HH^{\mathsf T}}9.
\]
Sa fonction caractéristique s’obtient en annulant les arguments du
second groupe de modes. Elle reste donc gaussienne et les théorèmes
précédents calculent tout son spectre, sans tronquer le registre d’occupation.
